% from ega4-15-terminal-fr.tex
\oldpage[IV-3]{253}
\phantomsection
\addcontentsline{toc}{section}{Table des matières}
\section*{TABLE DES MATIÈRES}

\begingroup
\small
\setlength{\parindent}{0pt}
\setlength{\parskip}{1.5pt}
\newcommand{\EGAIVthreeContentsLine}[3]{\textbf{#1}\quad #2\unskip\nobreak\hspace{0.4em}\nobreak\dotfill\nobreak\mbox{#3}\par}

\EGAIVthreeContentsLine{CHAPITRE IV.}{\textbf{Étude locale des schémas et des morphismes de schémas} \emph{(suite)}}{5}
\EGAIVthreeContentsLine{\S~8.}{Limites projectives de préschémas}{5}
\EGAIVthreeContentsLine{8.1.}{Introduction}{5}
\EGAIVthreeContentsLine{8.2.}{Limites projectives de préschémas}{7}
\EGAIVthreeContentsLine{8.3.}{Parties constructibles dans une limite projective de préschémas}{12}
\EGAIVthreeContentsLine{8.4.}{Critères d'irréductibilité et de connexion pour les limites projectives de préschémas}{17}
\EGAIVthreeContentsLine{8.5.}{Modules de présentation finie sur une limite projective de préschémas}{19}
\EGAIVthreeContentsLine{8.6.}{Sous-préschémas de présentation finie d'une limite projective de préschémas}{25}
\EGAIVthreeContentsLine{8.7.}{Critères pour qu'une limite projective de préschémas soit un préschéma réduit (resp. intègre)}{27}
\EGAIVthreeContentsLine{8.8.}{Préschémas de présentation finie sur une limite projective de préschémas}{28}
\EGAIVthreeContentsLine{8.9.}{Premières applications à l'élimination des hypothèses noethériennes}{34}
\EGAIVthreeContentsLine{8.10.}{Propriétés de permanence des morphismes par passage à la limite projective}{36}
\EGAIVthreeContentsLine{8.11.}{Application aux morphismes quasi-finis}{41}
\EGAIVthreeContentsLine{8.12.}{Nouvelle démonstration et généralisation du « Main Theorem » de Zariski}{43}
\EGAIVthreeContentsLine{8.13.}{Traduction en termes de pro-objets}{49}
\EGAIVthreeContentsLine{8.14.}{Caractérisation d'un préschéma localement de présentation finie sur un autre, en termes du foncteur qu'il représente}{52}

\EGAIVthreeContentsLine{\S~9.}{Propriétés constructibles}{54}
\EGAIVthreeContentsLine{9.1.}{Le principe de l'extension finie}{54}
\EGAIVthreeContentsLine{9.2.}{Propriétés constructibles et ind-constructibles}{56}
\EGAIVthreeContentsLine{9.3.}{Propriétés constructibles de morphismes de préschémas algébriques}{60}

\clearpage
